\chapter{NKContainers : ranger des données sans la bibliothèque standard}

Vous savez maintenant ce qu'est un template. \texttt{NKContainers} en est la
démonstration : cinquante-deux en-têtes qui donnent à Nkentseu ses tableaux, ses
chaînes, ses tables et ses arbres --- sans une ligne de bibliothèque standard.

Ce chapitre montre ceux dont vous vous servirez tous les jours, dit ce que
chacun coûte, et nomme les pièges que le dépôt a déjà payés.

\section{Pourquoi ne pas employer la bibliothèque standard}

Trois raisons, et aucune n'est idéologique.

\begin{center}
\begin{tabular}{@{}ll@{}}
\toprule
\textbf{Raison} & \textbf{Conséquence concrète} \\
\midrule
Plateformes contraintes & certaines n'ont pas de STL utilisable \\
Maîtrise de l'allocation & tout passe par NKMemory, donc se mesure \\
Comportement identique partout & la STL varie d'un compilateur à l'autre \\
\bottomrule
\end{tabular}
\end{center}

\begin{nkretenir}
La règle est \textbf{positive} : on n'écrit pas \texttt{std::}, et \emph{si la
primitive manque, on l'écrit au niveau où elle doit vivre} --- dans Foundation,
pas chez soi. Un utilitaire écrit localement devient le troisième exemplaire
d'une chose qui aurait dû être partagée.
\end{nkretenir}

\section{Le tableau dynamique : \texttt{NkVector<T>}}

C'est le conteneur que vous emploierez neuf fois sur dix.

\begin{nkcode}{Kernel/Foundation/NKContainers/src/NKContainers/Sequential/NkVector.h}
NkVector<int32> v;
v.PushBack(10);
v.PushBack(20);
v.EmplaceBack(30);            // construit SUR PLACE, pas de copie

const usize n = v.Size();     // combien d'elements
const usize c = v.Capacity(); // combien AVANT la prochaine reallocation
int32 premier = v[0];
v.Reserve(100);               // reserve d'un coup, evite N reallocations
v.Clear();                    // vide, GARDE la capacite
\end{nkcode}

\begin{center}
\begin{tabular}{@{}lll@{}}
\toprule
\textbf{Méthode} & \textbf{Effet} & \textbf{Coût} \\
\midrule
\texttt{PushBack(v)} & ajoute à la fin & amorti constant \\
\texttt{EmplaceBack(args\dots)} & construit à la fin & amorti constant \\
\texttt{operator[]} & accès direct & constant \\
\texttt{Insert(pos, v)} & insère au milieu & linéaire \\
\texttt{Erase(pos)} & retire & linéaire \\
\texttt{Find(v)} & cherche & linéaire \\
\texttt{Reserve(n)} & réserve & une allocation \\
\bottomrule
\end{tabular}
\end{center}

\begin{nkretenir}
\textbf{\texttt{Size} n'est pas \texttt{Capacity}.} La taille est le nombre
d'éléments présents ; la capacité, la place réservée. Quand la taille atteint la
capacité, le tableau réalloue --- il demande un bloc plus grand et recopie tout.

D'où \texttt{Reserve} : si vous savez que vous allez ajouter mille éléments, une
réservation remplace une dizaine de réallocations. Ce n'est pas une
micro-optimisation, c'est une allocation contre dix.
\end{nkretenir}

\begin{nkpiege}
\textbf{Toute réallocation invalide les pointeurs et les itérateurs.} Un
pointeur pris sur \texttt{v[0]}, puis un \texttt{PushBack} qui déclenche une
réallocation : le pointeur désigne l'ancien bloc, déjà libéré.

\begin{nkcode}{Exemple pedagogique -- FAUX}
int32 *p = &v[0];
v.PushBack(42);      // peut REALLOUER
*p = 7;              // pointeur pendant : on ecrit dans le vide
\end{nkcode}

C'est le \emph{pointeur pendant} du chapitre sur le C, sous sa forme la plus
courante. Prenez l'indice, pas l'adresse.
\end{nkpiege}

\begin{nkpiege}
\textbf{\texttt{Clear} ne libère pas la mémoire}, il remet la taille à zéro. Un
vecteur qui a contenu un million d'éléments occupe toujours autant après un
\texttt{Clear}. Pour rendre la mémoire : \texttt{ShrinkToFit}.
\end{nkpiege}

\section{Les chaînes : \texttt{NkString}}

\begin{nkcode}{Kernel/Foundation/NKContainers/src/NKContainers/String/NkString.h}
NkString s = "Bonjour";
s.Append(" tout le monde");
const char *brut = s.Data();          // pour les API qui veulent du C
const usize n = s.Size();
bool ok = s.StartsWith("Bon");
NkString morceau = s.SubStr(8);       // a partir de l'indice 8

NkString msg = NkString::Format("%d entite(s) en %0.2f ms", 42, 1.5f);
\end{nkcode}

\begin{nkpiege}
\textbf{\texttt{NkString::Format} est de la famille \texttt{printf}} : ses
marqueurs sont \texttt{\%d}, \texttt{\%s}, \texttt{\%f}. Le journal, lui, a
\emph{deux} familles --- \texttt{logger.Infof} en \texttt{printf}, et
\texttt{logger.Info} en marqueurs indexés \texttt{\{0\}}. Mélanger les deux
n'échoue pas : le marqueur s'affiche littéralement, le message sort faux et
ressemble à un message.
\end{nkpiege}

\begin{nkpiege}
\textbf{Une chaîne peut venir de nulle part.} \texttt{env::GetEnvVar} rend un
\texttt{const char *} qui vaut \texttt{nullptr} quand la variable n'existe pas.
Écrire \texttt{NkString v = env::GetEnvVar("X");} construit alors une chaîne
depuis un pointeur nul.

Ici le constructeur garde par \texttt{if (str)} et rend une chaîne vide --- mais
« absente » et « définie à vide » deviennent indistinguables. Si la différence
compte, testez le \texttt{const char *} \emph{avant} de le convertir.
\end{nkpiege}

\section{Les tables}

\begin{center}
\begin{tabular}{@{}lll@{}}
\toprule
\textbf{Conteneur} & \textbf{Ordonné ?} & \textbf{Recherche} \\
\midrule
\texttt{NkUnorderedMap<K,V>} & non & constante en moyenne \\
\texttt{NkMap<K,V>} & oui, par clé & logarithmique \\
\texttt{NkUnorderedSet<T>} & non & constante en moyenne \\
\texttt{NkSet<T>} & oui & logarithmique \\
\bottomrule
\end{tabular}
\end{center}

\begin{nkretenir}
Prenez la version \emph{non ordonnée} par défaut : elle est plus rapide. Ne
prenez l'ordonnée que si vous devez \textbf{parcourir dans l'ordre des clés} ---
c'est la seule chose qu'elle apporte, et elle se paie.
\end{nkretenir}

\section{Le reste, et à quoi ça sert}

\begin{center}
\begin{tabular}{@{}ll@{}}
\toprule
\textbf{Conteneur} & \textbf{Quand} \\
\midrule
\texttt{NkArray<T,N>} & taille fixe connue à la compilation, zéro allocation \\
\texttt{NkSpan<T>} & \emph{vue} sur des données qu'on ne possède pas \\
\texttt{NkOptional<T>} & « une valeur, ou rien » --- sans pointeur nul \\
\texttt{NkResult<T,E>} & « une valeur, ou une erreur » \\
\texttt{NkPair}, \texttt{NkTuple} & deux valeurs, ou \emph{n} \\
\texttt{NkVariant} & une valeur parmi plusieurs types \\
\texttt{NkFunction} & ranger un appel dans une variable \\
\texttt{NkRingBuffer<T>} & file de taille fixe qui tourne (audio, journaux) \\
\texttt{NkQueue}, \texttt{NkStack}, \texttt{NkDeque} & file, pile, file double \\
\texttt{NkPriorityQueue} & le plus urgent en premier \\
\texttt{NkList}, \texttt{NkDoubleList} & listes chaînées \\
\texttt{NkQuadTree}, \texttt{NkOctree} & partition de l'espace, 2D et 3D \\
\texttt{NkBTree}, \texttt{NkTrie}, \texttt{NkGraph} & arbres et graphes \\
\texttt{NkSort} & le tri du dépôt \\
\texttt{NkUTF8}, \texttt{NkUTF16}, \texttt{NkUTF32} & encodages de texte \\
\texttt{NkBase64} & encodage binaire vers texte \\
\bottomrule
\end{tabular}
\end{center}

\begin{nkretenir}
\textbf{\texttt{NkSpan<T>} ne possède rien} : c'est un pointeur et une longueur.
Il résout le piège du chapitre sur le C --- « un tableau ne connaît pas sa
taille » --- en transportant les deux ensemble. Mais il ne prolonge pas la vie de
ce qu'il désigne : un span sur un vecteur détruit est un pointeur pendant.
\end{nkretenir}

\begin{nkpiege}
\textbf{\texttt{NkOptional} n'est pas un pointeur nullable, et c'est mieux.}
\texttt{HasValue()} dit s'il y a quelque chose ; \texttt{Value()} le donne.
Appeler \texttt{Value()} sur un optional vide est une faute --- alors que
déréférencer un pointeur nul est un plantage. La différence : l'un se teste, et
la signature vous oblige à y penser.
\end{nkpiege}

\begin{nkbilan}
\texttt{NkVector} pour presque tout, \texttt{NkString} pour le texte,
\texttt{NkUnorderedMap} pour associer. \texttt{Size} n'est pas \texttt{Capacity},
et toute réallocation invalide les pointeurs. \texttt{Clear} ne rend pas la
mémoire. \texttt{NkSpan} transporte un tableau \emph{avec} sa taille, mais ne le
possède pas. \texttt{NkOptional} remplace le pointeur nullable par quelque chose
que la signature oblige à tester.
\end{nkbilan}

\begin{nkquestions}
\begin{enumerate}
\item Quelle différence entre \texttt{Size()} et \texttt{Capacity()} ?
\item Après \texttt{Clear()}, la mémoire est-elle rendue ? Comment la rendre ?
\item Pourquoi \texttt{Reserve} peut-il changer les performances d'un facteur
      important ?
\item Quand prendre \texttt{NkMap} plutôt que \texttt{NkUnorderedMap} ?
\item \texttt{NkSpan} possède-t-il ses données ? Qu'est-ce que cela implique ?
\end{enumerate}
\end{nkquestions}

\begin{nkpratique}
Écrivez une fonction qui reçoit un \texttt{NkVector<NkString>} et rend un
\texttt{NkUnorderedMap<NkString, int32>} comptant les occurrences de chaque
chaîne.

Puis mesurez : construisez un vecteur de 100\,000 chaînes, comptez avec et sans
\texttt{Reserve} sur la table, et comparez les durées avec \texttt{NkChrono}.
\end{nkpratique}

\begin{nkdemo}
Prouvez qu'une réallocation invalide les pointeurs.

Prenez un \texttt{NkVector<int32>}, réservez exactement 2 éléments, gardez
\texttt{\&v[0]}, puis ajoutez jusqu'à dépasser la capacité. Affichez la
\emph{valeur de l'adresse} de \texttt{\&v[0]} avant et après.

Deux adresses différentes prouvent la réallocation. \textbf{N'écrivez pas à
travers l'ancien pointeur} pour le « vérifier » : le programme pourrait ne rien
faire d'anormal, et vous concluriez à tort que c'est sûr.
\end{nkdemo}

\begin{nklogique}
\texttt{PushBack} est dit « à coût amorti constant » alors qu'une réallocation
recopie tous les éléments.

Supposez que la capacité \emph{double} à chaque réallocation. Comptez le nombre
total de recopies pour insérer $n$ éléments en partant d'une capacité de 1.
Montrez que ce total est proportionnel à $n$, et non à $n^2$ --- puis expliquez
en une phrase pourquoi doubler, plutôt qu'ajouter une place, change la nature du
coût.
\end{nklogique}
